% ============================================================
%  sections/chapter_05.tex
%  Section 5: Experiments
% ============================================================

\section{Experiments}
\label{sec:experiments}

We conduct four controlled numerical experiments, each designed to verify
a specific theoretical prediction from Sections~\ref{sec:background}--\ref{sec:norms}.
All experiments use a standard test image of size $256 \times 256 \times 3$
(RGB, values normalised to $[0, 1]$) unless otherwise stated.
Implementations follow the structure described in Section~\ref{sec:demo};
random seeds are fixed at 42 for reproducibility.

% ------------------------------------------------------------
\subsection{Experiment 1: Rank Truncation and Theoretical Error Bound}
\label{subsec:exp_rank}
% ------------------------------------------------------------

\paragraph{Hypothesis.}
The HOSVD error bound~\eqref{eq:hosvd_bound} states that
\begin{equation}
  \frobn{\tensor{X} - \hat{\tensor{X}}_{\mathrm{HOSVD}}}^2
  \;\leq\;
  B(r_1, r_2, r_3)
  \;\coloneqq\;
  \sum_{n=1}^{3} \sum_{i=r_n+1}^{I_n} \sigma_i^{(n)2}.
  \label{eq:exp1_bound}
\end{equation}
We expect the empirical error to track $B(r,r,r)$ as $r$ varies,
with HOOI strictly below HOSVD due to iterative refinement.

\paragraph{Setup.}
Tucker rank is swept as $r \in \{2, 4, 8, 16, 24, 32, 48, 64\}$
with the color axis fixed at $r_3 = 3$ (full rank).
For each $r$, we record:
\begin{itemize}
  \item Empirical HOSVD reconstruction error (left-hand side of~\eqref{eq:exp1_bound}),
  \item Theoretical bound $B(r, r, 3)$ (right-hand side),
  \item HOOI error after convergence,
  \item Compression ratio $\rho$~\eqref{eq:tucker_ratio},
  \item PSNR of the reconstructed image.
\end{itemize}

\paragraph{Expected results.}
\begin{itemize}
  \item Empirical HOSVD error $\leq B(r,r,3)$ for all $r$ (bound verified).
  \item HOOI error $\leq$ HOSVD error (refinement monotonically helps).
  \item Error decays at a rate governed by the singular value spectrum;
        rapid decay implies a low effective rank.
  \item PSNR crosses 30\,dB (visually acceptable) around $r \approx 20$--30
        for natural images.
\end{itemize}

\paragraph{Metrics reported.}
\begin{equation}
  \text{RelErr}
  \;=\; \frac{\frobn{\tensor{X} - \hat{\tensor{X}}}}{\frobn{\tensor{X}}},
  \qquad
  \text{PSNR}
  \;=\; 10 \log_{10}\!\left(\frac{1}{\mathrm{MSE}}\right)
  \;\text{[dB]},
  \qquad
  \rho
  \;=\; \frac{\prod_n I_n}{P_{\mathrm{Tucker}}}.
  \label{eq:exp1_metrics}
\end{equation}

% ------------------------------------------------------------
\subsection{Experiment 2: Norm Comparison under Structured Noise}
\label{subsec:exp_noise}
% ------------------------------------------------------------

\paragraph{Hypothesis.}
Under sparse corruption, $\mathcal{L}_1$-Tucker (IRLS) should outperform
$\mathcal{L}_F$-Tucker (HOOI), and the performance gap should grow with
corruption density (Section~\ref{subsec:norm_comparison}).
Under dense Gaussian noise, $\mathcal{L}_F$ should be optimal.

\paragraph{Setup.}
We construct three noise conditions:
\begin{enumerate}
  \item \textbf{Gaussian}: $e_{ijk} \sim \mathcal{N}(0, \sigma^2)$,
        $\sigma \in \{0.05, 0.10, 0.20\}$.
  \item \textbf{Sparse outliers}: a fraction $p \in \{0.05, 0.10, 0.20\}$
        of entries is replaced by $\mathrm{Uniform}(-1, 1)$ noise.
  \item \textbf{Mixed}: Gaussian ($\sigma = 0.05$) plus sparse corruption
        ($p = 0.10$).
\end{enumerate}
Tucker rank is fixed at $(r, r, 3)$ with $r = 20$.
Three methods are compared: HOOI ($\mathcal{L}_F$), IRLS ($\mathcal{L}_1$),
and Adam ($\mathcal{L}_H$, $\delta = 0.1$).

\paragraph{Expected results.}
\begin{itemize}
  \item Gaussian noise: HOOI achieves lowest RelErr; IRLS and Adam are
        close but slightly suboptimal.
  \item Sparse outliers: IRLS achieves lowest RelErr; advantage grows
        with $p$.
  \item Mixed: Huber (Adam) achieves best compromise.
  \item Error maps: HOOI spreads residual uniformly; IRLS concentrates
        residual on corrupted pixels.
\end{itemize}

% ------------------------------------------------------------
\subsection{Experiment 3: Convergence Analysis}
\label{subsec:exp_convergence}
% ------------------------------------------------------------

\paragraph{Hypothesis.}
HOOI should exhibit monotone decreasing error with fast initial convergence
(superlinear near the solution due to the SVD step).
Adam should converge more slowly but eventually reach similar quality
on the Huber objective.
IRLS outer iterations should converge in fewer than 20 steps.

\paragraph{Setup.}
All three methods are run for 100 iterations on the same tensor
($256 \times 256 \times 3$, rank $(20, 20, 3)$, no added noise).
We record at every iteration:
\begin{equation}
  e_t^{(F)} = \frobn{\tensor{X} - \hat{\tensor{X}}^{(t)}}, \qquad
  \Delta_t^{(n)} = \frobn{\mat{U}^{(n)(t)} - \mat{U}^{(n)(t-1)}},
  \label{eq:exp3_metrics}
\end{equation}
where $e_t^{(F)}$ measures reconstruction quality and
$\Delta_t^{(n)}$ measures factor matrix stability.

\paragraph{Expected results.}
\begin{itemize}
  \item HOOI: $e_t^{(F)}$ strictly non-increasing; $\Delta_t^{(n)} \to 0$.
  \item Adam: non-monotone $e_t^{(F)}$ but overall decreasing trend.
  \item IRLS: fast outer convergence (10--15 steps) when warm-started
        from HOOI solution.
\end{itemize}

% ------------------------------------------------------------
\subsection{Experiment 4: Initialization Sensitivity}
\label{subsec:exp_init}
% ------------------------------------------------------------

\paragraph{Hypothesis.}
Due to non-convexity, different initializations may converge to different
local minima.
HOSVD initialization should consistently outperform random initialization
and reduce variance across runs.

\paragraph{Setup.}
For Tucker rank $(20, 20, 3)$ on the clean image, we compare:
\begin{enumerate}
  \item \textbf{HOSVD} (deterministic, 1 run),
  \item \textbf{Random orthogonal} (5 independent runs, different seeds),
  \item \textbf{Random Gaussian} (5 independent runs).
\end{enumerate}
Each initialization is followed by 50 HOOI iterations.
We report the distribution (mean $\pm$ std) of final RelErr across runs.

\paragraph{Expected results.}
\begin{itemize}
  \item HOSVD initialization converges to a lower error than random.
  \item Random orthogonal has lower variance than random Gaussian.
  \item The gap between best and worst random run (max RelErr $-$ min RelErr)
        quantifies the severity of the local minima problem.
\end{itemize}

% ------------------------------------------------------------
\subsection{Summary of Results}
\label{subsec:exp_summary}
% ------------------------------------------------------------

Table~\ref{tab:exp_summary} consolidates the key quantitative results
across all four experiments.
Detailed figures (convergence curves, error maps, rank-error plots)
are presented alongside the image compression demonstration in
Section~\ref{sec:demo}.

\begin{table}[h]
  \centering
  \caption{Summary of experimental results.
           RelErr values are relative Frobenius errors.
           All results use Tucker rank $(20, 20, 3)$ unless stated.}
  \label{tab:exp_summary}
  \begin{tabular}{llcc}
    \toprule
    Experiment & Setting & Method & RelErr \\
    \midrule
    \multirow{2}{*}{Exp.\ 1 (rank sweep)}
      & $r = 20$  & HOSVD & --- \\
      & $r = 20$  & HOOI  & --- \\
    \midrule
    \multirow{3}{*}{Exp.\ 2 (noise, $p=0.10$)}
      & Sparse outliers & HOOI  & --- \\
      & Sparse outliers & IRLS  & --- \\
      & Sparse outliers & Adam  & --- \\
    \midrule
    \multirow{2}{*}{Exp.\ 3 (convergence)}
      & Iterations to $\Delta < 10^{-4}$ & HOOI & --- \\
      & Iterations to $\Delta < 10^{-4}$ & Adam & --- \\
    \midrule
    \multirow{2}{*}{Exp.\ 4 (init.)}
      & HOSVD init & HOOI & --- \\
      & Random init (mean $\pm$ std) & HOOI & --- \\
    \bottomrule
  \end{tabular}
\end{table}

\noindent
Dashes (---) are filled with experimental values upon code execution.
The table structure is provided here to make the theory--experiment
connection explicit before results are available.

% % ============================================================
% %  sections/chapter_05_ko.tex
% %  5장: 실험
% % ============================================================

% \section{실험}
% \label{sec:experiments}

% 본 장에서는 \ref{sec:background}--\ref{sec:norms}장의 이론적 예측을 각각 검증하도록
% 설계된 네 가지 통제된 수치 실험을 수행한다.
% 별도로 명시하지 않는 한, 모든 실험은 크기 $256 \times 256 \times 3$
% (RGB, 값 정규화 범위 $[0, 1]$)의 표준 테스트 이미지를 사용한다.
% 구현은 \ref{sec:demo}장에서 설명하는 구조를 따르며, 재현성을 위해
% 난수 시드는 42로 고정한다.

% % ------------------------------------------------------------
% \subsection{실험 1: 랭크 절단과 이론적 오차 상한}
% \label{subsec:exp_rank}
% % ------------------------------------------------------------

% \paragraph{가설.}
% HOSVD 오차 상한~\eqref{eq:hosvd_bound}에 의하면
% \begin{equation}
%   \frobn{\tensor{X} - \hat{\tensor{X}}_{\mathrm{HOSVD}}}^2
%   \;\leq\;
%   B(r_1, r_2, r_3)
%   \;\coloneqq\;
%   \sum_{n=1}^{3} \sum_{i=r_n+1}^{I_n} \sigma_i^{(n)2}.
%   \label{eq:exp1_bound}
% \end{equation}
% $r$이 변함에 따라 실험적 오차가 $B(r,r,r)$을 추적할 것으로 예상하며,
% HOOI는 반복 정제로 인해 HOSVD보다 엄밀히 낮은 오차를 보일 것으로 예상한다.

% \paragraph{설정.}
% Tucker 랭크를 $r \in \{2, 4, 8, 16, 24, 32, 48, 64\}$로 변화시키며,
% 색상 축은 $r_3 = 3$ (전체 랭크)으로 고정한다.
% 각 $r$에 대해 다음을 기록한다:
% \begin{itemize}
%   \item 실험적 HOSVD 재구성 오차 (식~\eqref{eq:exp1_bound} 좌변),
%   \item 이론적 상한 $B(r, r, 3)$ (우변),
%   \item 수렴 후 HOOI 오차,
%   \item 압축률 $\rho$~\eqref{eq:tucker_ratio},
%   \item 재구성 이미지의 PSNR.
% \end{itemize}

% \paragraph{예상 결과.}
% \begin{itemize}
%   \item 모든 $r$에서 실험적 HOSVD 오차 $\leq B(r,r,3)$ (상한 검증).
%   \item HOOI 오차 $\leq$ HOSVD 오차 (정제가 단조적으로 도움).
%   \item 오차 감소율은 특이값 스펙트럼에 의해 결정되며,
%         급격한 감소는 낮은 유효 랭크를 시사.
%   \item 자연 이미지에서 $r \approx 20$--30 근방에서 PSNR 30\,dB
%         (시각적으로 수용 가능) 초과.
% \end{itemize}

% \paragraph{보고 지표.}
% \begin{equation}
%   \text{RelErr}
%   \;=\; \frac{\frobn{\tensor{X} - \hat{\tensor{X}}}}{\frobn{\tensor{X}}},
%   \qquad
%   \text{PSNR}
%   \;=\; 10 \log_{10}\!\left(\frac{1}{\mathrm{MSE}}\right)
%   \;\text{[dB]},
%   \qquad
%   \rho
%   \;=\; \frac{\prod_n I_n}{P_{\mathrm{Tucker}}}.
%   \label{eq:exp1_metrics}
% \end{equation}

% % ------------------------------------------------------------
% \subsection{실험 2: 구조적 잡음 하에서의 노름 비교}
% \label{subsec:exp_noise}
% % ------------------------------------------------------------

% \paragraph{가설.}
% 희소 손상 하에서는 $\mathcal{L}_1$-Tucker (IRLS)가
% $\mathcal{L}_F$-Tucker (HOOI)를 능가해야 하며,
% 성능 격차는 손상 밀도에 따라 증가해야 한다
% (\ref{subsec:norm_comparison}절).
% 밀집 가우시안 잡음 하에서는 $\mathcal{L}_F$가 최적이어야 한다.

% \paragraph{설정.}
% 세 가지 잡음 조건을 구성한다:
% \begin{enumerate}
%   \item \textbf{가우시안}: $e_{ijk} \sim \mathcal{N}(0, \sigma^2)$,
%         $\sigma \in \{0.05, 0.10, 0.20\}$.
%   \item \textbf{희소 이상치}: 비율 $p \in \{0.05, 0.10, 0.20\}$의
%         원소를 $\mathrm{Uniform}(-1, 1)$ 잡음으로 대체.
%   \item \textbf{혼합}: 가우시안 ($\sigma = 0.05$) + 희소 손상 ($p = 0.10$).
% \end{enumerate}
% Tucker 랭크는 $r = 20$으로 고정하여 $(r, r, 3)$을 사용한다.
% 세 방법을 비교한다: HOOI ($\mathcal{L}_F$), IRLS ($\mathcal{L}_1$),
% Adam ($\mathcal{L}_H$, $\delta = 0.1$).

% \paragraph{예상 결과.}
% \begin{itemize}
%   \item 가우시안 잡음: HOOI가 최저 RelErr; IRLS와 Adam은 근접하나 약간 열등.
%   \item 희소 이상치: IRLS가 최저 RelErr; 이점이 $p$에 따라 증가.
%   \item 혼합: Huber (Adam)가 최선의 절충안 달성.
%   \item 오차 지도: HOOI는 잔차가 균등 분산; IRLS는 손상 픽셀에 집중.
% \end{itemize}

% % ------------------------------------------------------------
% \subsection{실험 3: 수렴 분석}
% \label{subsec:exp_convergence}
% % ------------------------------------------------------------

% \paragraph{가설.}
% HOOI는 단조 감소 오차를 보이며, SVD 단계로 인해 해 근방에서
% 초선형(superlinear) 수렴을 보일 것이다.
% Adam은 더 느리게 수렴하지만 Huber 목적 함수에서 유사한 품질에 도달할 것이다.
% IRLS 외부 반복은 20회 이내에 수렴할 것이다.

% \paragraph{설정.}
% 동일한 텐서($256 \times 256 \times 3$, 랭크 $(20, 20, 3)$,
% 추가 잡음 없음)에 세 방법을 100회 반복 실행한다.
% 매 반복마다 기록:
% \begin{equation}
%   e_t^{(F)} = \frobn{\tensor{X} - \hat{\tensor{X}}^{(t)}}, \qquad
%   \Delta_t^{(n)} = \frobn{\mat{U}^{(n)(t)} - \mat{U}^{(n)(t-1)}},
%   \label{eq:exp3_metrics}
% \end{equation}
% 여기서 $e_t^{(F)}$는 재구성 품질, $\Delta_t^{(n)}$는
% 인수 행렬 안정성을 측정한다.

% \paragraph{예상 결과.}
% \begin{itemize}
%   \item HOOI: $e_t^{(F)}$이 엄밀히 단조 비증가; $\Delta_t^{(n)} \to 0$.
%   \item Adam: $e_t^{(F)}$은 비단조이지만 전반적으로 감소 추세.
%   \item IRLS: HOOI 해로 웜 스타트 시 10--15회 외부 반복으로 빠른 수렴.
% \end{itemize}

% % ------------------------------------------------------------
% \subsection{실험 4: 초기화 민감도}
% \label{subsec:exp_init}
% % ------------------------------------------------------------

% \paragraph{가설.}
% 비볼록성으로 인해 서로 다른 초기화가 서로 다른 국소 최솟값으로
% 수렴할 수 있다.
% HOSVD 초기화는 무작위 초기화보다 일관되게 우수하며 런 간 분산을
% 줄여야 한다.

% \paragraph{설정.}
% 깨끗한 이미지에서 Tucker 랭크 $(20, 20, 3)$으로 비교한다:
% \begin{enumerate}
%   \item \textbf{HOSVD} (결정론적, 1회 실행),
%   \item \textbf{무작위 직교} (5회 독립 실행, 서로 다른 시드),
%   \item \textbf{무작위 가우시안} (5회 독립 실행).
% \end{enumerate}
% 각 초기화 후 50회 HOOI 반복을 수행한다.
% 런 전반에 걸쳐 최종 RelErr의 분포(평균 $\pm$ 표준편차)를 보고한다.

% \paragraph{예상 결과.}
% \begin{itemize}
%   \item HOSVD 초기화가 무작위보다 낮은 오차로 수렴.
%   \item 무작위 직교가 무작위 가우시안보다 낮은 분산.
%   \item 최선과 최악의 무작위 런 간 격차(최대 RelErr $-$ 최소 RelErr)가
%         국소 최솟값 문제의 심각도를 정량화.
% \end{itemize}

% % ------------------------------------------------------------
% \subsection{결과 요약}
% \label{subsec:exp_summary}
% % ------------------------------------------------------------

% 표~\ref{tab:exp_summary}은 네 실험의 핵심 정량적 결과를 통합한다.
% 상세 그림(수렴 곡선, 오차 지도, 랭크-오차 플롯)은
% \ref{sec:demo}장의 이미지 압축 데모와 함께 제시한다.

% \begin{table}[h]
%   \centering
%   \caption{실험 결과 요약.
%            RelErr은 상대 프로베니우스 오차.
%            별도 명시 없는 경우 Tucker 랭크 $(20, 20, 3)$ 사용.}
%   \label{tab:exp_summary}
%   \begin{tabular}{llcc}
%     \toprule
%     실험 & 설정 & 방법 & RelErr \\
%     \midrule
%     \multirow{2}{*}{실험 1 (랭크 스윕)}
%       & $r = 20$ & HOSVD & --- \\
%       & $r = 20$ & HOOI  & --- \\
%     \midrule
%     \multirow{3}{*}{실험 2 (잡음, $p=0.10$)}
%       & 희소 이상치 & HOOI  & --- \\
%       & 희소 이상치 & IRLS  & --- \\
%       & 희소 이상치 & Adam  & --- \\
%     \midrule
%     \multirow{2}{*}{실험 3 (수렴)}
%       & $\Delta < 10^{-4}$까지 반복 수 & HOOI & --- \\
%       & $\Delta < 10^{-4}$까지 반복 수 & Adam & --- \\
%     \midrule
%     \multirow{2}{*}{실험 4 (초기화)}
%       & HOSVD 초기화   & HOOI & --- \\
%       & 무작위 초기화 (평균 $\pm$ 표준편차) & HOOI & --- \\
%     \bottomrule
%   \end{tabular}
% \end{table}

% \noindent
% 대시(---)는 코드 실행 후 실험값으로 채워진다.
% 결과가 확보되기 전에도 이론--실험 연결을 명시적으로 드러내기 위해
% 표 구조를 미리 제시하였다.